\section{Solución analítica de problemas inversos lineales: guía por pseudoínverson}

Cuando el mapeo hacia adelante $\mathcal{A}$ es una función lineal (lo cual abarca una gran cantidad de aplicaciones prácticas como desenfocamiento, superresolución, restauración y reconstrucción TC), podemos obtener una expresión analítica para la función de verosimilitud.

\subsection{Configuración del modelo hacia adelante lineal}

Supongamos que $\mathcal{A}$ es lineal, es decir, $y = A x_0 + n$, donde $A$ es una matriz y $n \sim \mathcal{N}(0, \sigma_y^2 I)$ es el ruido de observación. En este caso:

$$
p(y | x_0) = \mathcal{N}(y; A x_0, \sigma_y^2 I)
$$

\begin{knowledgebox}{Matriz $A$ en problemas inversos lineales comunes}
\begin{itemize}
    \item \textbf{Desenfocamiento gaussiano}: $A$ es la matriz correspondiente al núcleo de convolución gaussiana.
    \item \textbf{Superresolución}: $A$ es la matriz de muestreo inferior (downsampling).
    \item \textbf{Restauración de imágenes}: $A$ es una matriz diagonal de máscara (las regiones conocidas son 1, las desconocidas son 0).
    \item \textbf{Reconstrucción TC/MRI}: $A$ es la matriz de proyección/muestreo de Fourier.
\end{itemize}
\end{knowledgebox}

\subsection{Derivación del prior gaussiano conjugado}

Utilizando las propiedades de conjugación de la distribución gaussiana, cuando el prior $p(x_0 | x_t)$ y la verosimilitud $p(y|x_0)$ son ambos distribuciones gaussianas, la distribución marginal $p(y | x_t)$ también es una distribución gaussiana. Específicamente:

\textbf{Prior} (después de aproximación gaussiana):
$$
p(x_0 | x_t) \approx \mathcal{N}(\hat{x}_0, r_t^2 I)
$$

\textbf{Verosimilitud}:
$$
p(y | x_0) = \mathcal{N}(A x_0, \sigma_y^2 I)
$$

\textbf{Marginal} (después de integración sobre $x_0$):
$$
p(y | x_t) = \int p(y | x_0) p(x_0 | x_t) \, dx_0 = \mathcal{N}(A \hat{x}_0, \; r_t^2 A A^\top + \sigma_y^2 I)
$$

\begin{importantbox}{Fórmula de guía por pseudoínverson (Pseudo-Inverse Guidance)}
Basado en la distribución gaussiana marginal anterior, el gradiente de verosimilitud puede calcularse analíticamente:
$$
\nabla_{x_t} \log p(y | x_t) \approx -\frac{1}{2} \nabla_{x_t} \left\| y - A \hat{x}_0 \right\|^2_{(r_t^2 A A^\top + \sigma_y^2 I)^{-1}}
$$
donde $\hat{x}_0$ es la estimación de eliminación de ruido calculada a partir de $x_t$ mediante la fórmula de Tweedie.

Cuando no hay ruido en la observación ($\sigma_y = 0$), la fórmula se simplifica a requerir el cálculo de la \textbf{pseudoínverson} de $A$, por lo que se denomina guía por pseudoínverson.
\end{importantbox}

\begin{warningbox}{La guía por pseudoínverson solo aplica a $\mathcal{A}$ lineal}
La derivación de esta fórmula analítica depende de la linealidad de $\mathcal{A}$. Para mapeos hacia adelante no lineales (como el renderizado que involucra redes neuronales o simulaciones físicas complejas), se deben utilizar métodos más generales (aunque aproximaciones más gruesas).
\end{warningbox}

\subsection{Resumen de la sección}

Para problemas inversos lineales, mediante las propiedades del prior gaussiano conjugado, podemos obtener una expresión analítica para el gradiente de verosimilitud. Este método requiere calcular la pseudoínverson de la matriz de mapeo hacia adelante, por lo que se denomina guía por pseudoínverson. Proporciona una solución teóricamente más rigurosa para tareas clásicas de procesamiento de imágenes como desenfocamiento, superresolución y restauración.
